%\input ../util/bookmacros
%\input ../util/docparams

\chapter{Introduction to Differential Equations}

\section{Motivational Problem}
The population of a species of bacteria, represented by $p(t)$, is 
the number of bacteria on a given surface. The evolution of the population is modeled by the 
assumption that its ``instantaneous'' growth rate
is proportional to the amount of bacteria currently present. The growth 
is impeded by a chemical on the surface that kills bacteria at a rate of $b$ bacteria per second.
These facts may be collected into an equation:
$$
\eqalign{
p'(t) & = r \, p(t) - b \cr
p(0) & = P \cr
}
$$
Here, $r$ is the proportionality coefficient and has units $1/{\rm second}$.
With this choice the units of the two sides of the equation match: the left
hand side is the limit,
$\lim_{n  \rightarrow \infty} {p(t+h_n) - p(t) \over h_n}$ (where $h_n \rightarrow 0$),
which has units of bacteria per second; on the right hand side, $r \, p(t)$ has
units $1/{\rm second} \times {\rm bacteria}$ -- bacteria per second again --
matching $b$, which is given in bacteria per second.


This equation is called a differential equation which you can think of as an
infinitesimal difference equation. To see why, chop time into small steps of
length $h$ and let $x_n = p(n\,h)$. Over one step the population changes by
approximately $h$ times its rate of change:
$$
\eqalign{
p(t + h) & \approx p(t) + h\,\big(r\, p(t) - b\big) \cr
x_{n+1} & = x_n + h\,(r\, x_n - b) \cr
}
$$
The second line is a difference equation of exactly the kind we solved in the
last chapter; the differential equation is what remains as $h \rightarrow 0$.
This is the {\it approximate, then take a limit\/} move run in reverse.
How do we solve this and what does a solution mean? The solution is 
a {\it function\/} $p$ which satisfies the above equations. We will find an explicit formula for the solution.

\section{Basic Solution Strategy}
Let's start with a simpler differential equation. In fact, let's start with 
the simplest differential equation that we can think of.
$$
\eqalignno{
x'(t) & = 0 & (1) \cr
x(0) & = P \cr
}
$$
To solve this equation we must find a {\it function\/} such that 
its derivative is $0$. That is, the tangent at each point is flat.
We know $x$ must be a constant function. Since the initial 
value is $P$, the function must be $x(t) = P$.
Another way to solve this equation is to use the \index{Fundamental Theorem 
of Calculus}. We do this by integrating out the derivative. 
Integrating (1) from $0$ to $t$ gives
$$
\eqalign{
\int_{0}^t x'(s) \, ds &= \int_{0}^t 0 \, ds  \cr
x(t) - x(0) & = 0 
}
$$
Therefore, $x(t) = x(0)$, and since $x(0) = P$, we have 
$x(t) = P$.

As with \index{difference equation}s, there is no technique for writing 
a formula for the solution of a general \index{differential equation}.
Again, as with difference equations, there are limited classes 
of equations for which we can write down a solution. However, unlike difference 
equations there is a question of existence and uniqueness of 
a solution. It turns out there are \index{differential equation}s 
which don't have solutions, and others for which the solution 
is not unique. When we say that a solution doesn't exist, we don't 
just mean that we can't find a formula to express the solution; 
we mean that there are some equations for which there is no 
function which satisfies the differential equation. And, there are differential equations which have more than 
one solution.%
\footnote{\kern 0.5pt \raise 0.5ex \hbox{\dag}}{A discussion of these issues is beyond the scope of this book.}

The way to find formulas for the solutions to these problems is 
to use the \index{Fundamental Theorem of Calculus} to integrate away the 
derivative. This is analogous to the way we solved difference equations:
we used a summation operator to get rid of the \index{telescoping difference}s. 
That is, we used the \index{Fundamental Theorem of Difference Calculus} to find solutions.

\subsection{Solution of the Motivational Problem}
Let's look at an equation that is closer to our motivational problem.
Consider the simpler problem:
$$
\eqalignno{
p'(t) & =  r\, p(t) & (2) \cr
p(0) & = P & \cr
}
$$
If we integrate both sides we can use the \index{Fundamental Theorem of Calculus} 
to integrate away the derivative on the left hand side of the equation.
However, we don't know what the integral of the right hand side is since it
involves the unknown function, $p$.
We need a way to combine all the $p$'s together and view them as the 
derivative of some function of $p$. Then, the Fundamental Theorem 
of Calculus can be used to eliminate the derivative. We then will have an 
equation which does not have derivatives, and therefore, we can use algebra 
to find $p$. To do this
divide both sides by $p$ -- assuming $p(t) > 0$, which is reasonable for a
population; we now have all the $p$'s (the unknowns) on one side
of the equation. (The {\it integrating factor\/} method we use below for the
full problem avoids this division entirely.)
$$
\eqalign{
{p'(t) \over p(t)} & =  r  \cr
p(0) & = P \cr
}
$$
The left hand side is the derivative of $\ln(p)$.\footnote{\kern 0.5pt \raise 0.5ex \hbox{\ddag}}
{This follows from the chain rule.} Integrating 
both sides from $0$ to $t$ we find:
$$
\eqalign{
\ln(p(t)) - \ln(p(0)) & =  r(t - 0) \cr
}
$$
Or,
$$
\eqalign{
\ln(p(t)) & = \ln(p(0)) + r\,t \cr
}
$$
The inverse function of $\ln(x)$ is $e^x$; therefore, exponentiating 
both sides of the equation gives us
$$
p(t) = e^{\ln(p(0)) + r\,t}  = e^{\ln(p(0))} e^{r\, t} = p(0) \,e^{r \,t} = P \, e^{r \,t}
$$
To solve our motivational problem, we try to do the same thing; that is,
arrange the $p$'s on one side in such a way that they are the derivative 
of some expression involving $p$. Then, use the Fundamental Theorem of Calculus to 
integrate away the derivative. We then have an ordinary equation for which 
we solve for $p$. 

To do this we employ a technique analogous to what was 
done for the motivational problem in the chapter on difference equations.
First move the $p$ on the right hand side to the left hand side -- grouping 
all the $p$'s together. Next, 
multiply both sides by some (as yet unknown) function $y(t)$. Our motivational problem becomes:
$$
\eqalignno{
 y(t)  p'(t) - r y(t) \,p(t) & = -b\,y(t) & (3)
}
$$
The goal is to choose $y$ so that the left hand side is the derivative of the product, $(y(t) p(t))$.
The derivative of this product is
$$
(y(t) p(t))' = (p(t) y(t))' = p'(t) y(t) + p(t) y'(t) = y(t) p'(t) + y'(t) p(t)
$$
To match the left hand side of (3) we need $y$ to satisfy $y'(t) = -r \, y(t)$.
But this is an equation we know how to solve, the solution is $y(t) = C\, e^{-r\, t}$, for some constant $C$.
Since there are no further restrictions on $y$, we take $C$ to be 1.
Such a multiplier function $y$ is called an {\it \index{integrating factor}\/} --
the continuous twin of the {\it summing factor\/} we met in the chapter on
difference equations.
So, equation (3) becomes
$$
\big(e^{-r \, t}p(t)\big)' = -b\, e^{-r \, t}
$$
Integrating both sides of this equation from $0$ to $t$ gives:%
\footnote{\kern 0.5pt \raise 0.5ex \hbox{\dag}}%
{We're using integration to collapse the infinitesimal telescoping differences; that is, 
we use the Fundamental Theorem of Calculus to get a simple difference. 
Notice we used a different integration variable, $s$, to avoid confusion with $t$.}
$$
\int_0^t \big(e^{-r \, s}p(s)\big)' \, ds = \int_0^t -b\, e^{-r \, s} \, ds
$$
It is not clear what the anti-derivative of $e^{-r \, s}$ is however. We do know that the derivative 
of $e^{a\, x}$ is $a\, e^{a \, x}$ from the subsection ``Derivatives of Log
and Exponential Functions'' in the chapter on Differential Calculus. From this we can guess that the anti-derivative of $e^{a\, x}$ is ${e^{a \, x} \over a}$;
which can be verified by taking its derivative.

Now, the integral above can be evaluated to give
$$
e^{-r\, t}p(t) - e^{ -r \, 0} p(0) = {b \over r} e^{-r \, t} - {b \over r}
$$
Solving for $p$ we obtain (since $e^0 = 1$ and $p(0) = P$)
$$
p(t) = P\, e^{r \, t} +  {b \over r}(1 - e^{r \, t})
$$
This may be written:
$$
p(t) = {b \over r} + (P - {b \over r}) e^{r \, t}
$$
This solution has the same structure as the loan solution of the last chapter.
The constant piece, ${b \over r}$, is an {\it equilibrium\/} solution of the
full equation (if $P = {b \over r}$, then $p'(t) = 0$ and the population never
moves); the exponential piece, $(P - {b \over r})\, e^{r \, t}$, solves the
no-harvest equation $p'(t) = r \, p(t)$. Because the equation is
{\it linear\/}, the two pieces simply add -- the {\it exploit linearity\/}
move ({\it \index{superposition}\/}) again.

Notice that if $P - {b \over r} < 0$ then $p(t)$ will eventually become zero. Let us assume that this 
is the case and find the time to extinction. That is, find $t^{*}$ such that $p(t^{*}) = 0$. The
equation to solve for $t^*$ (put $t^*$ on one side and everything else on the other) is
$$
p(t^*) = 0 = {b \over r} + (P - {b \over r}) e^{r \, t^*}
$$
After some algebra this becomes
$$
e^{r\, t^*} = - {{ b \over r} \over {P - {b \over r}}} = {{ b \over r} \over {{b \over r} - P}}
$$
Taking the $\log_e$ (that is, $\ln$) of both sides and solving for $t^*$ gives
$$
t^* = {1 \over r} \ln\bigg( {{ b \over r} \over {{b \over r} - P}}\bigg) =  {1 \over r} \ln\bigg( {b  \over {b - r \, P}}\bigg) 
$$
So, in the case that $P < {b \over r}$ (which is the same as $r \, P < b$) the population of bacteria dies out at this
$t^*$. That makes sense since initially the rate of change of the population is $p'(0) = r \, p(0) - b = r \, P - b < 0$;
so the population starts out decreasing. A smaller population makes
$p'(t) = r \, p(t) - b$ even more negative, and the decline feeds on itself.
So, once the population starts to shrink it becomes a death spiral for the bacteria.

The whole family of solutions can be seen in one picture.
\epsfigure{eps/bacteria.eps}{The family of solutions
$p(t) = {b \over r} + (P - {b \over r})\, e^{r\,t}$ for several starting
populations $P$. The dashed line is the equilibrium $p = b/r$: populations
starting above it grow without bound, populations starting below it spiral
down to extinction.}

It is interesting to note that the formula for $p(t)$ becomes negative for $t > t^*$. Obviously, this
doesn't make physical sense -- the model no longer holds after the critical time $t^*$.
Notice also that our model predicts that there are a non integral number of bacteria at almost any given time.
The model could be corrected by removing the fractional part of the predicted population size by either truncating or
rounding the result. One could think of this extra step as part of an improved model. 

Keep in mind that models are always
an approximation to the truth. For instance, although the model that we have may work over short intervals of time, there is
no population that satisfies this equation for too long because as the population rises other factors in the environment
inhibit growth.

\subsection{Radioactive Decay}
We apply this technique to a problem in physics.

A radioactive substance decays at a rate proportional to the amount present. If one starts with $Q$
pounds of the substance, in 1,000 years there will be $Q/4$ pounds. When will there be $\frac{1}{2}Q$ pounds?

The equation is similar to the previous problem:
$$
\eqalign{
a'(t) & = -r\, a(t) \cr
a(0) & = Q \cr
}
$$
Here, $a(t)$ is the amount in pounds at time $t$ in years; and, $r$ is taken to be a positive constant, whose 
units must be 1 / years for the equation to make sense. We don't know what $r$ is but we know that 
in 1,000 years there will be a quarter of $Q$ left. Solving the equation gives $a(t) = Q\, e^{-r t}$. 
When $t$ is 1,000 the equation reads
$a(1000) = Q\, e^{-1000 r} = \frac{1}{4}\,Q$. We may solve this equation for $r$ by taking the $\log_e$ (that is, $\ln$)
of the last equation. 
This gives us $\ln(Q) - 1000\, r = \ln(\frac{1}{4}Q) = \ln(1) - \ln(4)  + \ln(Q) = \ln(Q)
-\ln(4)$ (since $\ln(1) = 0$). 
Or, $-1000 \, r = -\ln(4)$.
So, $r = \frac{\ln(4)}{1000}$.
To find the time $t$ when $a(t) = \frac{1}{2}\, Q$, we need to solve:
$a(t) =Q\, e^{-r t} = \frac{1}{2}\, Q$ for $t$ (the only unknown in the equation is $t$ as we now know $r$).
Taking the $\ln$  of both sides gives us
$\ln(Q) -r\, t = \ln(\frac{1}{2}Q) = \ln(1) - \ln(2) + \ln(Q) = \ln(Q) - \ln(2) $. 
So $t = \ln(2) / r = 1000  \ln(2) / \ln(4) = 500$. 

That is, the time it takes
the initial amount to decay to $\frac{1}{2}Q$ pounds is 500 years.
Notice that the answer is independent of the initial amount $Q$. In other words, the ``half-life'' of
this \index{radioactive material} is an intrinsic attribute of this substance.

\exercise{(Newton's Law of Cooling) A cup of coffee at temperature $T(t)$ sits
in a room at constant temperature $A$. It cools at a rate proportional to the
difference between its temperature and the room's:
$T'(t) = -k\,\bigl(T(t) - A\bigr)$, with $T(0) = T_0$ and $k$ a positive
constant. Solve for $T(t)$ using an integrating factor. What does your
solution predict as $t$ grows large?}

\exercise{A radioactive substance has a half-life of $100$ years. How long
until only $10$\% of the original amount remains? (First use the half-life to
find $r$, as in the subsection above.)}

\section{Summary}
The strategy of this chapter is the one we have been playing since first-year
algebra: {\it isolate and invert\/}. We isolate the unknown function under the
derivative operator $d/dt$ -- multiplying by an {\it integrating factor\/}
when needed to make the left hand side a single derivative -- and then invert
$d/dt$ by applying its inverse, integration, in the form of the
\index{Fundamental Theorem of Calculus}. This is precisely what we did in the
last chapter with $\Delta$, summation, and the Fundamental Theorem of
Difference Calculus. The whole book culminates in the following dictionary:

\bigskip
\centerline{\vbox{\halign{\quad\hfil#\hfil\quad&\quad\hfil#\hfil\quad\cr
{\bf Discrete}&{\bf Continuous}\cr
\noalign{\smallskip\hrule\smallskip}
$x_{n+1} = a\, x_n$&$p'(t) = r\, p(t)$\cr
geometric growth, $a^n$&exponential growth, $e^{r\,t}$\cr
$\Delta$&$d/dt$\cr
summation&integration\cr
FTDC&FTC\cr
summing factor&integrating factor\cr
}}}
\bigskip

The two columns are tied together by the correspondence $a = 1 + r$: for small
$r$ we have $(1+r)^n \approx e^{r\,n}$, which is why the solution of the loan
problem and the solution of the population problem look so alike.

\exercise{Make the correspondence concrete: for the loan of the last chapter
the monthly rate was $r = 0.005$ and the term was $N = 60$ months. Compare
$(1.005)^{60}$ with $e^{0.005 \times 60} = e^{0.3}$. How close are they?}

And just as the Fibonacci sequence hinted at a second order discrete story,
there is a continuous second order story -- $x'' + x = 0$, the equation of the
swinging pendulum. Its full treatment awaits a course on differential
equations, but you can already verify its heart -- and cash the promise made
when the trig derivatives were introduced:

\exercise{Using the derivatives of sine and cosine from the chapter on
Differential Calculus, check that $x(t) = \sin(t)$ and $x(t) = \cos(t)$ both
satisfy $x'' + x = 0$. There is the oscillation the trig functions were
built for. Then check that {\it any\/} combination
$A\sin(t) + B\cos(t)$ is also a solution -- the {\it exploit linearity\/}
move (superposition) once more, just as it was for the two Fibonacci
solutions $\lambda_1^n$ and $\lambda_2^n$.}

\section{What's Left}
A \index{differential equation} is effectively any equation which involves a \index{derivative}. Obviously, they can
be much more complicated than the ones we looked at. A careful study of the subject classifies 
differential equations into various categories and then provides recipes for solutions in some cases. 
The main solution technique is to
``integrate away'' the derivatives; that is, use the \index{Fundamental Theorem of Calculus}.
